\par\medskip
\noindent\textup{(D)}
至少在一个适用 (C) 的情形中, $\delta_x$ 在 $\ell$ 进上同调中不是挠元.

设 $S_1$ 为 $\operatorname{Spec}(\mathbb Z)$ 在 $(p)$ 处的严格 Hensel 化, $V\subset\mathbb P^r/S_1$ 为 $S_1$ 上纯相对维数为 $n+1$ 的光滑射影簇. 假设 $V/S_1$ 的特殊纤维 $V_p$ 存在超平面截面的一个 Lefschetz 线束(XVII, 2.5). 此束由与 $\mathbb P^r/S_1$ 对偶的射影空间 $\check{\mathbb P}^{,r}/S_1$ 的特殊纤维中的一条直线 $d'$ 定义. 将 $d'$ 提升为 $\check{\mathbb P}^{,r}/S_1$ 中的一条直线 $d$, 并设 $X\to d$ 为由 $d$ 参数化的 $V$ 的超平面截面束(XVIII, 3.1). $d$ 中的例外轨迹 $\Delta$ 在 $S_1$ 上平展(XVII, 6). 因而在 $\Delta_p\subset d$ 的某点处作平展局部化后, 得到 (C) 类型的情形.

为检验 $\delta_x$ 非挠, 将标量从 $S_1$ 扩张到 $\mathbb C$ 后检验即可. 取 $V=\mathbb P^{n+1}$, 其射影嵌入由 $\mathcal O(N)$ 定义. 当 $N\ge2$ 时可应用 XVII, 2.5(在此特例中对每个 $N$ 均成立). 然后使用以下引理.

\subsubsection*{引理 3.3.7.}
\underline{对 $\mathbb P^{n+1}(\mathbb C)$ 中次数 $N$ 充分大的超曲面的 Lefschetz 线束($N\ge3$ 即可), $\mathbb Q_\ell$-上同调中的消灭圈非零.}

此引理由 Lefschetz 线束理论得出, 例如由 XVIII, 6.6.1 得出. 此处使用 XVIII 中基于“Picard--Lefschetz 公式”的结果并非循环论证, 因为我们已有 (A). 也可改用更直接且纯超越的变体 XIX, no.~4, 它仅依赖于 XIV, 3.2.

\medskip
\noindent\textup{(E)}
我们在剩余特征为 $p$ 的情形对 $(X/S,x\in X_s)$ 证明 3.3.5.

设 $X_0/S_0$ 如 (D). 由 1.3.2(i), 存在态射 $g:S\to S_0$, 使
\[
g^*X_0=X_0\times_{S_0}S
\]
在 $x_0$ 处的 Hensel 化 $S$-同构于 $X$ 在 $x$ 处的 Hensel 化. 因此只需对 $g^*X_0$ 证明 3.3.5. 对 $X_0/S_0$ 的 (C) 作基变换, 在 $H^n(X_{0\bar\eta})$ 中得到
\[
\sigma(a)=a-(-1)^m
\varepsilon_{b(x)}(\sigma)(a,\delta_x)\delta_x
\]
其中 $\sigma$ 属于惯性群. 设 $\pi$ 为一致化元, 将此公式与 (2.2.5.9) 比较, 其中 $\varepsilon=\varepsilon_\pi$. 在 $H^n(X_{0\bar\eta})$ 中得到
\[
\lambda_X\varepsilon_\pi(\sigma)(a,\delta_x)\delta_x
=-(-1)^m\varepsilon_{b(x)}(\sigma)(a,\delta_x)\delta_x
\]
由于 $\delta_x$ 非挠, 可在此等式中约去 $(a,\delta_x)\delta_x$. 将得到的 $\lambda_X$ 值代入公式 (2.2.5.9), 恰好得到 3.3.5.

\medskip
\noindent\textup{(F) 一般情形下 3.3.6 的证明.}

3.3.6 的一般情形归结为 3.3.5 的一般情形, 正如 (B) 和 (C) 归结为 (A).